\section{階層線形モデル}
\subsection{授業内容}
\subsubsection{科目の中でのこのコマの位置づけ}
ここまで線形モデル，一般化線形モデルの例を見てきたが，\citet{久保2012}の例にならって一般化線形混合モデル，階層ベイズモデルへとモデルを展開させていく。

一般化線形混合モデルについては，Withinデザインですでに対応しており，正規分布以外の確率分布を使うことで一般化可能である。ここに切片や傾きなど，係数の方に分布が混ぜ合わせられることで階層化されたモデルとなる。ネストされたデータの具体例として，反復測定と大規模調査の二種類を取り上げ，また階層モデルの設計図を書いてから分析コードを書く手順を確認する。
\subsubsection{コマ主題細目}
\begin{description}
	\item[一般化線形混合モデル] Withinデザインの分析モデルを再確認するところから考える。このモデルは個体差のような個人ごとに変わる要因が含まれており，ここで変量効果と固定効果の違いを考える。また従属変数が正規分布でないモデルにすることで，一般化線形モデルと考えることができる。
	\item[ネストされたデータ] すでに反復測定データの場合が該当するが，データが階層性を持っているネストされたデータの例として，プロ野球データや大規模調査の例を考える。これらに対して，階層化しない分析とする分析とで解釈が異なる例を挙げ，階層的なデータに対して適切な分析が必要であることを理解する。
	\item[階層線形モデル] 階層化されたモデルを数式的に理解する。名称として，レベル1/2の効果，個人/集団レベルの効果と呼ばれることもあることを確認する。またモデルの設計図をかき，それをコードに起こすことで分析ができることを確認する。個別の回帰直線を引く場合に比べて，縮小が起こっていることをモデル比較を通じて確認する。
\end{description}
\subsubsection{キーワード}
\begin{itemize}
	\item 混合モデル
	\item ネストされたデータ
	\item 固定効果
	\item 変量効果
	\item 階層モデル
\end{itemize}
\subsection{授業情報}
\paragraph{コマの展開方法}
講義/遠隔可/演習
\subsubsection{予習・復習課題}
\paragraph{予習}反復測定デザインの分析例を確認しておく。そこで分布が混合されていることを確認しておく。
\paragraph{復習}階層線形モデルが応用できるようなデータ例を身の回りから考えてみるとよい。その上で，データがどのような背景から生成されているかの設計図を書き，設計図からコードに起こすという手順を一歩ずつ確認しておこう。